% !TeX root = ../main.tex

Die Grafik~\ref{fig:d-rateExec} zeigt die Erfolgs- und Fehlerraten, sowie die Ausführungszeiten der drei Fehlerklassen
für die betrachteten Delay Werte.
Die Erfolgsrate von Deadlocks orientiert sich an den 50\%.
Bei einem Delay von 0s liegt sie bei 43.65\%, steigt zunächst an und fällt bei einem Delay von 0.05s wieder
auf 43.69\%.
Die höchste Erfolgsrate wird mit 48.16\% bei einem Delay von 0.1s erreicht.
Dabei haben Deadlocks im Vergleich die höchsten Ausführungszeiten.
Die Ausführungszeit erfolgreicher Transaktionen liegt über die betrachteten Delay-Werte hinweg bei ungefähr 5000ms.
Bei nicht erfolgreichen Transaktionen ist nach einem Delay von 0.005s nur noch eine geringe Veränderung der
Ausführungszeit zu erkennen.
Serialization Failure haben für alle Delay Werte außer 0s eine Erfolgsrate über 50\%.
Sie sind außerdem die einzigen die eine höhere Erfolgsrate als Fehlerrate haben.
Die höchste Erfolgsrate liegt bei 70\% bei einem Delay von 0.05s.
Serialization Failures weisen gleichzeitig die niedrigsten Ausführungszeiten auf.
Sowohl bei erfolgreichen als auch bei nicht erfolgreichen Transaktionen steigt die Ausführungszeit mit zunehmendem
Delay.
Bei erfolgreichen Transaktionen fällt der Anstieg geringer aus.
Die Ausführungszeit nicht erfolgreiche Transaktionen steigt 134ms bei 0s auf 222ms bei 0.005s und bleibt für 0.01s auf
diesem Niveau.
Bei 0.05s steigt sie auf 360ms und erreicht bei 0.1s den höchsten Wert mit 508ms.
Die Ausführungszeit erfolgreicher Transaktionen liegt zwischen 150ms und 160ms für die Delay-Werte von 0s bis 0.01s
und steigt anschließend auf 215ms bei 0.1s.
Die Erfolgsrate von Lock Timeouts ist über alle Delay Werte die niedrigste.
Diese liegt bei einem Delay von 0s bei 10\% und steigt über die betrachteten Werte auf 18\% bei 0.1s.
Für Lock Timeouts liegt die Ausführungszeit einer erfolgreichen Transaktion durchschnittlich bei 1000ms bis 1400ms.
Sie sinkt bei 0.005s leicht, steigt für 0.01s und 0.05s weiter an und sinkt für 0.1s wieder ab.
Die Ausführungszeit erfolgreicher Transaktionen liegt dabei über der von nicht erfolgreichen Transaktionen.
Letztere steigt dagegen nahezu linear von 383ms bei 0s auf 785ms bei 0.1s.
Für den Delay von 0s ist zu beachten, dass in \texttt{d0} das Baseline-Experiment enthalten ist.
Bei Serialization Failures liegt die Erfolgsrate beispielsweise für die \emph{immediate Retry}-Strategie bei 57\%.
Der Wert für 0s umfasst zusätzlich das \emph{No Retry}-Experiment.

% !TeX root = ../main.tex

\begin{figure}[ht]
	\centering
	\begin{minipage}{0.45\textwidth}
		\centering
		\caption*{(a) Erfolgs- und Fehlerraten}
		% !Tex root = ../main.tex

\begin{tikzpicture}
	\begin{axis}[
		width=0.8\textwidth,
		height=0.8\textwidth,
		ylabel={Rate (\%)},
		xlabel={Delay},
		xtick={1,2,3,4,5},
		xticklabels={d0,d0.005,d0.01,d0.05,d0.1},
		xticklabel style={font=\scriptsize},
		ymin=0,
		ymax=100,
		ytick={0,20,40,60,80,100},
		yticklabel style={font=\small},
		scaled y ticks=false,
		legend style={
			at={(0.5,-0.30)},
			anchor=north,
			legend columns=2,
			align=left,
			legend cell align=left
		},
	]

		% Deadlock
		\addplot[
			color=successcolor,
			solid,
			thick,
			mark=*,
			mark size=1pt
		] coordinates {
			(1,52.136)
			(2,46.632)
			(3,46.600)
			(4,43.688)
			(5,48.160)
		};

		\addplot[
			color=failedcolor,
			solid,
			thick,
			mark=*,
			mark size=1pt
		] coordinates {
			(1,47.864)
			(2,53.368)
			(3,53.400)
			(4,56.312)
			(5,51.840)
		};


		% Serialization Failure
		\addplot[
			color=successcolor,
			dashed,
			thick,
			mark=diamond*,
			mark size=2pt
		] coordinates {
			(1,57.570667)
			(2,57.408000)
			(3,56.206857)
			(4,70.056000)
			(5,69.536000)
		};


		\addplot[
			color=failedcolor,
			dashed,
			thick,
			mark=diamond*,
			mark size=2pt
		] coordinates {
			(1,42.429333)
			(2,42.592000)
			(3,43.793143)
			(4,29.944000)
			(5,30.464000)
		};


		% Lock Timeout
		\addplot[
			color=successcolor,
			dotted,
			thick,
			mark=triangle*,
			mark size=2pt
		] coordinates {
			(1,12.773333)
			(2,12.648000)
			(3,14.197714)
			(4,17.888000)
			(5,18.856000)
		};


		\addplot[
			color=failedcolor,
			dotted,
			thick,
			mark=triangle*,
			mark size=2pt
		] coordinates {
			(1,87.226667)
			(2,87.352000)
			(3,85.802286)
			(4,82.112000)
			(5,81.144000)
		};


	\end{axis}
\end{tikzpicture}
	\end{minipage}
	\begin{minipage}{0.45\textwidth}
		\centering
		\caption*{(b) Ausführungszeiten}
		% !TeX root = ../main.tex

\begin{tikzpicture}
	\begin{axis}[
		width=0.8\textwidth,
		height=0.8\textwidth,
		ylabel={Execution Time (ms)},
		xlabel={Delay},
		xtick={1,2,3,4,5},
		xticklabels={d0,d0.005,d0.01,d0.05,d0.1},
		xticklabel style={font=\scriptsize},
		ymax=15000,
		ymode=log,
		ytick={100,200,500,1000,2000,5000,10000,15000},
		yticklabels={100,200,500,1000,2000,5000,10000,15000},
		yticklabel style={font=\small},
		scaled y ticks=false,
		legend style={
			at={(0.5,-0.30)},
			anchor=north,
			legend columns=2,
			legend cell align=left
		},
	]
		% Deadlock
		\addplot[
			color=successcolor,
			solid,
			thick,
			mark=*,
			mark size=1pt
		] coordinates {
			(1,4558.899840)
			(2,4422.962086)
			(3,4607.303470)
			(4,4712.113166)
			(5,4799.195183)
		};

		\addplot[
			color=failedcolor,
			solid,
			thick,
			mark=*,
			mark size=1pt
		] coordinates {
			(1,8559.730115)
			(2,12272.201019)
			(3,11290.482012)
			(4,11986.756642)
			(5,13306.593827)
		};

		% Serialization Failure
		\addplot[
			color=successcolor,
			dashed,
			thick,
			mark=diamond*,
			mark size=2pt
		] coordinates {
			(1,149.116139)
			(2,156.745959)
			(3,158.162115)
			(4,206.039625)
			(5,215.703981)
		};

		\addplot[
			color=failedcolor,
			dashed,
			thick,
			mark=diamond*,
			mark size=2pt
		] coordinates {
			(1,134.121373)
			(2,222.656837)
			(3,222.712597)
			(4,359.798023)
			(5,507.839286)
		};

		% Timeout
		\addplot[
			color=successcolor,
			dotted,
			thick,
			mark=triangle*,
			mark size=2pt
		] coordinates {
			(1,1135.150159)
			(2,1080.852625)
			(3,1273.949610)
			(4,1362.268336)
			(5,1209.853627)
		};

		\addplot[
			color=failedcolor,
			dotted,
			thick,
			mark=triangle*,
			mark size=2pt
		] coordinates {
			(1,382.882614)
			(2,486.299844)
			(3,532.888168)
			(4,647.833106)
			(5,784.960268)
		};

	\end{axis}
\end{tikzpicture}
	\end{minipage}
	\pgfplotslegendfromname{gemeinsameLegende}
	\caption{Raten und durchschnittliche Ausführungszeit für einzelne Delay Werte}
	\label{fig:d-rateExec}
\end{figure}

Die Tabelle~\ref{tab:d-retryDistribution} zeigt die Verteilung der benötigten Retries an.
Ein Delay von 0s entspricht der Strategie \emph{immediate Retry}.
Diese wird in Abschnitt~\ref{subsec:s} genauer betrachtet.
Für den Delay von 0.01s werden die beiden Experimente \texttt{Experiment2} und \texttt{Experiment3}
aufgeteilt, damit mit man einen Vergleich zwischen dem Einfluss von Delay ziehen kann.
Auch die Quantile in Tabelle~\ref{tab:d-quantiles} zeigen überwiegend die bereits bei den durchschnittlichen
Ausführungszeiten beschriebenen Verläufe.
Eine Auffälligkeit zeigt sich bei Serialization Failures: Im $p_{99.9}$-Quantil liegt die Ausführungsdauer bei einem
Delay von 0.005s deutlich über dem $p_{99}$-Quantil.
Ein ähnlicher Abstand ist bei einem Delay von 0.05s zu beobachten.

Der Retry Overhead in Abbildung~\ref{fig:d-overhead} weist keine wesentlichen Abweichungen von den bereits
beschriebenen Ausführungszeiten auf.
Der Overhead von Deadlocks bleibt über die Delay-Werte weitgehend konstant, während sich bei Serialization Failures ein
Anstieg mit zunehmendem Delay zeigt.
Eine Abweichung tritt bei Lock Timeouts bei einem Delay von 0.1s auf, hier liegt der Overhead erfolgreicher
Transaktionen unter dem von nicht erfolgreichen Transaktionen.
Der Durchsatz in Abbildung~\ref{fig:d-throughput} liefert ebenfalls keine wesentlich neuen Beobachtungen.
Bei Deadlocks bleibt er über die Delay-Werte weitgehend konstant und weist bei 0.05s einen kurzen Rückgang auf.
Bei Serialization Failures und Lock Timeouts steigt der Durchsatz mit zunehmendem Delay.
Die weiteren Experimente sind der Vollständigkeit halber dargestellt, enthalten jedoch keine Variation des Delays.